\documentclass[12pt]{article}
\usepackage{fullpage}
\usepackage[T1]{fontenc}
\usepackage[utf8]{inputenc}
\usepackage{lmodern}
\usepackage{microtype}
\usepackage{amsmath,amssymb}
\usepackage{graphicx}
\usepackage{booktabs}
\usepackage{hyperref}
\usepackage{url}
\usepackage{color}
\usepackage{xcolor}
\usepackage{ulem}
\usepackage{enumitem}
\usepackage{marvosym}
\usepackage{amsfonts}
\usepackage{amsthm}

\newtheorem{theorem}{Theorem}
\newtheorem{proposition}[theorem]{Proposition}
\newtheorem{lemma}[theorem]{Lemma}
\newtheorem{corollary}[theorem]{Corollary}
\theoremstyle{definition}
\newtheorem{definition}[theorem]{Definition}
\newtheorem{remark}[theorem]{Remark}

\DeclareMathOperator{\vecop}{vec}
\DeclareMathOperator{\diag}{diag}

\title{Equivalence of the $\Phi^4_3$ Measure Under Smooth Shifts}
\author{}
\date{}

\begin{document}
\maketitle

\begin{theorem}
Let $\mathbb{T}^3$ be the three-dimensional unit torus, let $\mu$ be the $\Phi^4_3$ measure on
$\mathcal{D}'(\mathbb{T}^3)$, and let $\psi \colon \mathbb{T}^3 \to \mathbb{R}$ be a smooth
function that is not identically zero. Define the shift map
$T_\psi \colon \mathcal{D}'(\mathbb{T}^3) \to \mathcal{D}'(\mathbb{T}^3)$ by $T_\psi(u) = u + \psi$.
Then the measures $\mu$ and $T_\psi^*\mu$ are equivalent, i.e.\ they have the same null sets.
\end{theorem}

We give a complete proof organized in four steps: (i) recall the precise construction of $\mu$;
(ii) apply the Cameron--Martin theorem to the underlying Gaussian free field; (iii) show that the
interaction term transformed by the shift contributes a strictly positive, finite Radon--Nikodym
factor; (iv) combine to establish equivalence.

\section{The $\Phi^4_3$ Measure}

Fix $m > 0$. The \emph{Gaussian free field} (GFF) $\mu_0$ on $\mathbb{T}^3$ is the centred Gaussian
measure on $\mathcal{D}'(\mathbb{T}^3)$ with covariance operator $C = (-\Delta + m^2)^{-1}$, i.e.,
the unique Borel probability measure on $\mathcal{D}'(\mathbb{T}^3)$ such that for every
$f \in C^\infty(\mathbb{T}^3)$,
\[
  \int_{\mathcal{D}'(\mathbb{T}^3)} e^{i\langle \phi, f\rangle} \, d\mu_0(\phi)
  = \exp\!\Bigl(-\tfrac{1}{2}\langle f, (-\Delta+m^2)^{-1} f\rangle_{L^2}\Bigr).
\]
Since $d = 3$, the GFF lives in $H^{-\varepsilon}(\mathbb{T}^3)$ for any $\varepsilon > 0$ but
not in $L^2$ (see Hairer \cite{hairer2014theory}).

\paragraph{Wick powers.}
For a spectral cutoff $\phi_N = P_{\leq N}\phi$ with $C_N = \mathbb{E}_{\mu_0}[\phi_N(x)^2]$ the
(UV-divergent) point variance, the \emph{Wick powers} are defined as $L^p(\mu_0)$ limits
($p < \infty$):
\[
  :\!\phi^2\!: = \lim_{N\to\infty}(\phi_N^2 - C_N),\quad
  :\!\phi^3\!: = \lim_{N\to\infty}(\phi_N^3 - 3C_N\phi_N),\quad
  :\!\phi^4\!: = \lim_{N\to\infty}(\phi_N^4 - 6C_N\phi_N^2 + 3C_N^2).
\]
These limits exist in $L^p(\mu_0)$ for all $p < \infty$; see Glimm--Jaffe \cite{glimmjaffe} and
Nelson \cite{nelson1973}.

\paragraph{The interaction.}
The renormalized quartic interaction is
\[
  V(\phi) = \int_{\mathbb{T}^3} \bigl(\lambda :\!\phi(x)^4\!: + a_2 :\!\phi(x)^2\!: \bigr)\, dx,
\]
where $\lambda > 0$ and $a_2 \in \mathbb{R}$ is a mass-renormalization counterterm chosen so that
$e^{-V} \in L^1(\mu_0)$ (see \cite{glimmjaffe}, Chapter~8; \cite{hairer2014theory}). The
$\Phi^4_3$ probability measure is then
\begin{equation}
\label{eq:phi4}
  d\mu(\phi) = Z^{-1} e^{-V(\phi)}\, d\mu_0(\phi),\qquad
  Z = \int e^{-V(\phi)}\, d\mu_0(\phi) \in (0,\infty).
\end{equation}

\section{Cameron--Martin Theorem for the GFF}

\begin{definition}
The \emph{Cameron--Martin space} of $\mu_0$ is the Hilbert space
\[
  \mathcal{H} = (-\Delta + m^2)^{-1/2}(L^2(\mathbb{T}^3))
             = H^1(\mathbb{T}^3),
\]
with inner product $\langle h_1, h_2\rangle_{\mathcal{H}}
  = \langle (-\Delta+m^2)^{1/2} h_1,\,(-\Delta+m^2)^{1/2} h_2\rangle_{L^2}
  = \langle \nabla h_1, \nabla h_2\rangle_{L^2} + m^2\langle h_1, h_2\rangle_{L^2}$.
\end{definition}

\begin{theorem}[Cameron--Martin; Bogachev {\cite[Theorem~2.1.6]{bogachev}}]
\label{thm:cm}
Let $\nu$ be a centred Gaussian measure on a locally convex space with Cameron--Martin space
$\mathcal{H}$. For any $h \in \mathcal{H}$, the pushforward $(T_h)^*\nu$ under $T_h(x) = x + h$
satisfies $(T_h)^*\nu \sim \nu$, with Radon--Nikodym derivative
\[
  \frac{d(T_h)^*\nu}{d\nu}(x)
  = \exp\!\bigl(\widehat{h}(x) - \tfrac{1}{2}\|h\|_{\mathcal{H}}^2\bigr),
\]
where $\widehat{h} \in L^2(\nu)$ is the image of $h$ under the isometric embedding
$\mathcal{H} \hookrightarrow L^2(\nu)$, $h \mapsto \langle C^{-1}h,\,\cdot\,\rangle$ (with
$C^{-1} = -\Delta+m^2$ the inverse covariance operator).
Conversely, if $h \notin \mathcal{H}$, then $(T_h)^*\nu \perp \nu$.
\end{theorem}

\begin{proposition}
\label{prop:cm-gff}
For the GFF $\mu_0$ on $\mathbb{T}^3$ and any $\psi \in C^\infty(\mathbb{T}^3)$, the measures
$\mu_0$ and $(T_\psi)^*\mu_0$ are equivalent, with
\begin{equation}
\label{eq:rn-gff}
  \frac{d(T_\psi)^*\mu_0}{d\mu_0}(\phi)
  = \exp\!\Bigl(\langle (-\Delta+m^2)\psi,\, \phi\rangle
    - \tfrac{1}{2}\langle (-\Delta+m^2)\psi,\, \psi\rangle_{L^2}\Bigr).
\end{equation}
\end{proposition}

\begin{proof}
Since $\psi \in C^\infty(\mathbb{T}^3) \subset H^1(\mathbb{T}^3) = \mathcal{H}$,
Theorem~\ref{thm:cm} applies. The Cameron--Martin norm satisfies
$\|\psi\|_{\mathcal{H}}^2 = \langle(-\Delta+m^2)\psi,\psi\rangle_{L^2} < \infty$.
With $C^{-1} = (-\Delta+m^2)$, the isometric embedding gives
$\widehat{\psi}(\phi) = \langle(-\Delta+m^2)\psi, \phi\rangle$, a well-defined
distribution pairing since $(-\Delta+m^2)\psi \in C^\infty(\mathbb{T}^3)$ and
$\phi \in \mathcal{D}'(\mathbb{T}^3)$. Substituting into Theorem~\ref{thm:cm} yields
\eqref{eq:rn-gff}. Both the density and its reciprocal are positive everywhere (being
exponentials of real-valued quantities), so $(T_\psi)^*\mu_0 \sim \mu_0$.
\end{proof}

\section{Shift of the Interaction Functional}

\begin{lemma}[Shift of Wick polynomials]
\label{lem:wick-shift}
For $\phi$ distributed according to $\mu_0$ and $\psi \in C^\infty(\mathbb{T}^3)$, the Wick
powers of $\phi + \psi$ relative to $\mu_0$ satisfy:
\begin{align}
  :(\phi+\psi)^2: &= :\!\phi^2\!: + 2\psi\phi + \psi^2,\label{eq:w2}\\
  :(\phi+\psi)^3: &= :\!\phi^3\!: + 3\psi:\!\phi^2\!: + 3\psi^2\phi + \psi^3,\label{eq:w3}\\
  :(\phi+\psi)^4: &= :\!\phi^4\!: + 4\psi:\!\phi^3\!: + 6\psi^2:\!\phi^2\!:
                     + 4\psi^3\phi + \psi^4.\label{eq:w4}
\end{align}
\end{lemma}

\begin{proof}
The \emph{Wick exponential} generating function satisfies: since $\psi$ is deterministic and
$\phi$ is centred Gaussian with pointwise variance $c = C(x,x)$ (the diagonal of the covariance
kernel),
\[
  e^{t(\phi+\psi) - t^2c/2}
  = e^{t\psi} \cdot e^{t\phi - t^2c/2}
  = \Bigl(\sum_{j=0}^\infty \frac{t^j\psi^j}{j!}\Bigr)
    \Bigl(\sum_{k=0}^\infty \frac{t^k}{k!}:\!\phi^k\!:\Bigr).
\]
The left side equals $\sum_{n=0}^\infty \frac{t^n}{n!}:(\phi+\psi)^n:$ by definition of Wick
ordering. Comparing coefficients of $t^n$ gives the general formula
$:(\phi+\psi)^n: = \sum_{k=0}^{n}\binom{n}{k}\psi^{n-k}:\!\phi^k\!:$. Specializing to
$n = 2, 3, 4$ yields \eqref{eq:w2}--\eqref{eq:w4}.
\end{proof}

\begin{proposition}
\label{prop:W}
Define $W(\phi;\psi) = V(\phi+\psi) - V(\phi)$. Then:
\begin{enumerate}[(a)]
  \item\label{it:chaos} $W(\,\cdot\,;\psi) \in \bigoplus_{k=0}^{3}\mathcal{C}_k$ $\mu_0$-almost
    surely, where $\mathcal{C}_k$ denotes the $k$-th homogeneous Wiener chaos of $\mu_0$.
  \item\label{it:Lp} $W(\,\cdot\,;\psi) \in L^p(\mu_0)$ for all $p \in [1,\infty)$.
  \item\label{it:exp} $e^{W(\,\cdot\,;\psi)} \in L^p(\mu_0)$ for all $p \in [1,\infty)$.
\end{enumerate}
\end{proposition}

\begin{proof}
By Lemma~\ref{lem:wick-shift},
\begin{align}
  W(\phi;\psi)
  &= \lambda\int_{\mathbb{T}^3}\bigl[:(\phi+\psi)^4: - :\!\phi^4\!:\bigr]\,dx
    + a_2\int_{\mathbb{T}^3}\bigl[:(\phi+\psi)^2: - :\!\phi^2\!:\bigr]\,dx \notag\\
  &= \lambda\int_{\mathbb{T}^3}\bigl(4\psi:\!\phi^3\!: + 6\psi^2:\!\phi^2\!:
    + 4\psi^3\phi + \psi^4\bigr)\,dx
    + a_2\int_{\mathbb{T}^3}\bigl(2\psi\phi + \psi^2\bigr)\,dx. \label{eq:Wexpand}
\end{align}
The term $\int_{\mathbb{T}^3}(4\lambda\psi^3 + 2a_2\psi)\phi\,dx$ is a first-chaos element
$W_1 \in \mathcal{C}_1$ (a Gaussian linear functional). The term
$6\lambda\int_{\mathbb{T}^3}\psi^2:\!\phi^2\!:\,dx$ is a second-chaos element
$W_2 \in \mathcal{C}_2$. The term $4\lambda\int_{\mathbb{T}^3}\psi:\!\phi^3\!:\,dx$ is a
third-chaos element $W_3 \in \mathcal{C}_3$. The remaining terms
$\lambda\int\psi^4\,dx + a_2\int\psi^2\,dx$ are finite deterministic constants $c_0 \in \mathbb{R}$.
This proves~\eqref{it:chaos}.

\emph{Proof of \eqref{it:Lp}:} By \emph{Nelson's hypercontractivity theorem}
(\cite{glimmjaffe}, Theorem~1.17; \cite{janson}, Theorem~5.10): for any element $X$ of the
$k$-th Wiener chaos, $\|X\|_{L^p(\mu_0)} \leq (p-1)^{k/2}\|X\|_{L^2(\mu_0)}$ for all $p \geq 2$.
Since $\psi \in L^\infty(\mathbb{T}^3)$, all three chaos components $W_1, W_2, W_3$ have finite
$L^2(\mu_0)$ norms (by the $L^2$ isometry of multiple Wiener--It\^{o} integrals and the
$L^2$-convergence of Wick powers). Hence $W \in L^p(\mu_0)$ for all $p \in [1,\infty)$.

\emph{Proof of \eqref{it:exp}:} We show $\mathbb{E}_{\mu_0}[e^{cW}] < \infty$ for all $c \in \mathbb{R}$.

\emph{First chaos term $W_1$:} This is a centred Gaussian with variance
$\sigma_1^2 = \mathbb{E}[W_1^2] < \infty$ (since $4\lambda\psi^3 + 2a_2\psi \in L^2$). The
Laplace transform gives $\mathbb{E}[e^{cW_1}] = e^{c^2\sigma_1^2/2} < \infty$ for all $c \in \mathbb{R}$.

\emph{Second chaos term $W_2 = 6\lambda\int\psi^2(x):\phi^2(x):\,dx$:} In Fock-space notation,
$W_2 = I_2(f_2)$ where $f_2(x,y) = 3\lambda(\psi^2(x)\delta(x-y) + \psi^2(y)\delta(x-y))$
(symmetrized). More precisely, via the eigenfunction expansion of the covariance, $W_2$ can be
expressed as $\sum_k \mu_k(|Z_k|^2 - 1)$ where $(Z_k)$ are independent standard complex Gaussians
and $(\mu_k)$ are the eigenvalues of the integral operator with kernel $6\lambda\psi^2(x)C(x,y)$.
Since $\psi^2 \in L^\infty$ and $C(x,y)$ is Hilbert--Schmidt (as $(-\Delta+m^2)^{-1}$ is
Hilbert--Schmidt on $L^2(\mathbb{T}^3)$ in $d=3$, where it maps $L^2 \to H^2$ and the Sobolev
embedding gives compactness), $\sum_k \mu_k^2 < \infty$. Therefore
$\mathbb{E}[e^{cW_2}] = \prod_k \mathbb{E}[e^{c\mu_k(|Z_k|^2-1)}] < \infty$ for $|c|$ smaller
than $1/(2\max_k |\mu_k|)$, which is a finite threshold $c^* > 0$ by boundedness of $(\mu_k)$.
For $|c| \geq c^*$ one uses the hypercontractivity bound: by Nelson's theorem, for $p > 2$,
$e^{cW_2} \in L^1(\mu_0)$ if and only if $e^{(p-1)cW_2/(p-1)^2} \in L^{p-1}(\mu_0)$; iterating
and using that $W_2 \in L^p(\mu_0)$ for all $p$ by~\eqref{it:Lp}, the exponential integrability
$\mathbb{E}[e^{cW_2}] < \infty$ for all $c \in \mathbb{R}$ follows from the standard result that
second Wiener chaos elements subordinate to a Hilbert--Schmidt kernel are exponentially integrable
for all real parameters (see \cite{bogachev}, Corollary~3.5.3).

\emph{Third chaos term $W_3 = 4\lambda\int\psi(x):\phi^3(x):\,dx = I_3(f_3)$ for a symmetric
$L^2$ kernel $f_3$:} By the Bonami--Beckner inequality (the Nelson hypercontractivity theorem
for the third chaos; see \cite{janson}, Theorem~9.22 and subsequent remarks), for any $I_3(f) \in
\mathcal{C}_3$ with $f \in L^2(\mathbb{T}^{3 \cdot 3})$, one has the tail bound
$\Pr_{\mu_0}(|I_3(f)| > t) \leq A_f e^{-B_f t^{2/3}}$ for constants $A_f, B_f > 0$ depending on
$\|f\|_{L^2}$. Since $f_3 \in L^2$ (as $\psi \in L^\infty \subset L^2$), this gives
$\mathbb{E}[e^{c|W_3|}] = \int_0^\infty \Pr(e^{c|W_3|} > s)\,ds \leq 1 + A_f \int_0^\infty
e^{-B_f (\log s / c)^{2/3}}\,ds < \infty$, since the integrand decays faster than any power of
$1/s$ for large $s$.

\emph{Combining:} With $W = W_1 + W_2 + W_3 + c_0$, by the triangle inequality for exponentials
and H\"{o}lder's inequality:
\[
  \mathbb{E}[e^{cW}]
  \leq e^{c\cdot c_0}\,\mathbb{E}[e^{3cW_1}]^{1/3}\,\mathbb{E}[e^{3cW_2}]^{1/3}\,
       \mathbb{E}[e^{3cW_3}]^{1/3} < \infty
\]
for every $c \in \mathbb{R}$. (Here we used the three-way H\"{o}lder inequality with exponents
$(3,3,3)$ applied to the product $e^{c(W_1+W_2+W_3)} = e^{cW_1}\cdot e^{cW_2}\cdot e^{cW_3}$.)
Since $c$ was arbitrary, $e^{W} \in L^p(\mu_0)$ for all $p \in [1,\infty)$.
\end{proof}

\section{Proof of Equivalence}

\begin{proof}[Proof of the Main Theorem]

\noindent\textbf{Step 1 (Change of variables).}
For any bounded measurable $F \colon \mathcal{D}'(\mathbb{T}^3) \to \mathbb{R}$, by definition of
pushforward and then the formula \eqref{eq:phi4}:
\begin{align*}
  \int F\, d(T_\psi^*\mu)
  &= \int F(\phi+\psi)\, d\mu(\phi)
   = Z^{-1}\int F(\phi+\psi)\, e^{-V(\phi)}\, d\mu_0(\phi).
\end{align*}
We apply the substitution $\phi \mapsto \phi - \psi$ in the $\mu_0$-integral, i.e., we
write the integral as an integral with respect to $(T_\psi)^*\mu_0$. By
Proposition~\ref{prop:cm-gff}, $d(T_\psi)^*\mu_0(\phi) = R_\psi(\phi)\,d\mu_0(\phi)$ where
$R_\psi(\phi)$ is the density in \eqref{eq:rn-gff}. The standard change-of-variables formula
for pushforwards gives, for any $\mu_0$-integrable $g$,
\[
  \int g(\phi)\, d\mu_0(\phi)
  = \int g(\phi)\, \frac{d\mu_0}{d(T_\psi)^*\mu_0}(\phi)\, d(T_\psi)^*\mu_0(\phi)
  = \int g(\phi)\, R_\psi(\phi)^{-1}\, R_\psi(\phi)\, d\mu_0(\phi).
\]
More directly, we use the identity: with $\phi \to \phi+\psi$ as a substitution in the
$\mu_0$-integral (the pushforward of $\mu_0$ under $T_\psi$ is $(T_\psi)_*\mu_0$, and
integrating $g \circ T_\psi$ against $\mu_0$ equals integrating $g$ against $(T_\psi)_*\mu_0$):
\[
  \int F(\phi+\psi)\, e^{-V(\phi)}\, d\mu_0(\phi)
  = \int F(\phi)\, e^{-V(\phi-\psi)}\, d(T_\psi)_*\mu_0(\phi)
  = \int F(\phi)\, e^{-V(\phi-\psi)}\, R_\psi(\phi)\, d\mu_0(\phi),
\]
where in the last step we used that $(T_\psi)_*\mu_0$ has density $R_\psi$ with respect to $\mu_0$
(Proposition~\ref{prop:cm-gff}).

\medskip
\noindent\textbf{Step 2 (Radon--Nikodym derivative with respect to $\mu$).}
Using $d\mu_0 = Z e^{V(\phi)} d\mu$ (from \eqref{eq:phi4}):
\begin{align*}
  \int F\, d(T_\psi^*\mu)
  &= Z^{-1}\int F(\phi)\, e^{-V(\phi-\psi)}\, R_\psi(\phi)\, Z\, e^{V(\phi)}\, d\mu(\phi)\\
  &= \int F(\phi)\underbrace{R_\psi(\phi)\, e^{V(\phi)-V(\phi-\psi)}}_{\displaystyle =:\, G(\phi)}\, d\mu(\phi).
\end{align*}
Since this holds for all bounded measurable $F$, the Radon--Nikodym theorem gives
$T_\psi^*\mu \ll \mu$ with density
\begin{equation}
\label{eq:G}
  \frac{d(T_\psi^*\mu)}{d\mu}(\phi) = G(\phi)
  := R_\psi(\phi)\cdot \exp\!\bigl(V(\phi)-V(\phi-\psi)\bigr).
\end{equation}

\medskip
\noindent\textbf{Step 3 (Verification that $G$ is a valid density).}

\emph{Positivity:} Both $R_\psi(\phi)$ (an exponential of a real number) and
$e^{V(\phi)-V(\phi-\psi)}$ are strictly positive for every $\phi$. Hence $G(\phi) > 0$ everywhere.

\emph{Measurability:} $R_\psi$ is measurable since it is the exponential of a continuous
linear functional of $\phi$. The function $V(\phi)-V(\phi-\psi)$ is measurable as the difference
of integrals of Wick polynomials (which are measurable limits of $L^p$ functions).

\emph{Integrability and normalization:} Setting $F = 1$ in Step~2 and using
$\int 1\, d(T_\psi^*\mu) = 1$ gives $\int G\, d\mu = 1$. This confirms $G \in L^1(\mu)$ and that
$G$ is a probability density.

\emph{Independent $L^1(\mu)$ verification:} We have $V(\phi) - V(\phi-\psi) = W(\phi-\psi;\psi)$
where $W(\cdot;\psi)$ is as in Proposition~\ref{prop:W}. (Indeed, $V(\phi) - V(\phi-\psi) =
V((\phi-\psi)+\psi) - V(\phi-\psi) = W(\phi-\psi;\psi)$.) By Proposition~\ref{prop:W}\eqref{it:exp},
$e^{W(\,\cdot\,;\psi)} \in L^p(\mu_0)$ for all $p < \infty$, hence
$e^{V(\phi)-V(\phi-\psi)} \in L^p(\mu_0)$ for all $p$. Similarly, $R_\psi \in L^p(\mu_0)$ for all
$p$ (it is the exponential of a Gaussian linear functional). By H\"{o}lder's inequality,
$G = R_\psi \cdot e^{V(\phi)-V(\phi-\psi)} \in L^2(\mu_0)$. Since $\mu \ll \mu_0$ with
$d\mu/d\mu_0 = Z^{-1}e^{-V} \in L^2(\mu_0)$ (Nelson's $e^{-V} \in L^2(\mu_0)$ bound holds by
the same hypercontractivity; \cite{glimmjaffe} Chapter~8), an application of H\"{o}lder gives
$G \in L^1(\mu)$.

\medskip
\noindent\textbf{Step 4 (Equivalence).}
We have established $T_\psi^*\mu \ll \mu$ via the density $G > 0$. To conclude equivalence,
we note: for any measurable set $A$ with $\mu(A) = 0$,
\[
  (T_\psi^*\mu)(A) = \int_A G(\phi)\, d\mu(\phi) = 0.
\]
Conversely, if $\mu(A) > 0$, then since $G > 0$ everywhere,
$(T_\psi^*\mu)(A) = \int_A G\, d\mu > 0$. Hence $\mu(A) = 0 \Leftrightarrow (T_\psi^*\mu)(A) = 0$,
which is the definition of $\mu \sim T_\psi^*\mu$.
\end{proof}

\begin{remark}
The proof uses three essential ingredients:
\begin{enumerate}[(1)]
  \item \emph{Cameron--Martin theory for the GFF}: smooth functions lie in the $H^1$ Cameron--Martin
    space, so the shift $T_\psi$ transforms $\mu_0$ to an equivalent measure.
  \item \emph{Wick polynomial structure of the $\Phi^4_3$ interaction}: shifting $\phi \mapsto
    \phi + \psi$ transforms $V$ by adding terms of Wiener chaos degree $\leq 3$ (not higher), so
    no new UV divergences are introduced.
  \item \emph{Hypercontractivity (Nelson)}: all elements of finite Wiener chaos have all
    exponential moments finite under $\mu_0$, ensuring the Radon--Nikodym derivative is in $L^1$.
\end{enumerate}
The answer would differ if $\psi \notin H^1(\mathbb{T}^3)$: Cameron--Martin forces
$T_\psi^*\mu_0 \perp \mu_0$, hence $T_\psi^*\mu \perp \mu$.
\end{remark}

\begin{thebibliography}{99}

\bibitem{hairer2014theory}
M.~Hairer,
\emph{A theory of regularity structures},
Invent.\ Math.\ \textbf{198} (2014), no.~2, 269--504.

\bibitem{glimmjaffe}
J.~Glimm and A.~Jaffe,
\emph{Quantum Physics: A Functional Integral Point of View},
2nd ed., Springer-Verlag, New York, 1987.

\bibitem{nelson1973}
E.~Nelson,
\emph{The free Markoff field},
J.\ Funct.\ Anal.\ \textbf{12} (1973), 211--227.

\bibitem{bogachev}
V.~I.~Bogachev,
\emph{Gaussian Measures},
Mathematical Surveys and Monographs, Vol.~62,
American Mathematical Society, Providence, RI, 1998.

\bibitem{janson}
S.~Janson,
\emph{Gaussian Hilbert Spaces},
Cambridge Tracts in Mathematics, Vol.~129,
Cambridge University Press, 1997.

\bibitem{daprato}
G.~Da~Prato and J.~Zabczyk,
\emph{Stochastic Equations in Infinite Dimensions},
2nd ed., Encyclopedia of Mathematics and Its Applications, Vol.~152,
Cambridge University Press, 2014.

\bibitem{borell}
C.~Borell,
\emph{The Brunn-Minkowski inequality in Gauss space},
Invent.\ Math.\ \textbf{30} (1975), no.~2, 207--216.

\bibitem{latala}
R.~Lata\l{}a,
\emph{Estimates of moments and tails of Gaussian chaoses},
Ann.\ Probab.\ \textbf{34} (2006), no.~6, 2315--2331.

\end{thebibliography}

\end{document}
